\documentclass{article}


\usepackage{PRIMEarxiv}

\usepackage[utf8]{inputenc} % allow utf-8 input
\usepackage[T1]{fontenc}    % use 8-bit T1 fonts
\usepackage{hyperref}       % hyperlinks
\usepackage{url}            % simple URL typesetting
\usepackage{booktabs}       % professional-quality tables
\usepackage{amsfonts}       % blackboard math symbols
\usepackage{nicefrac}       % compact symbols for 1/2, etc.
\usepackage{microtype}      % microtypography
\usepackage{lipsum}
\usepackage{fancyhdr}       % header
\usepackage{graphicx}       % graphics
\graphicspath{{media/}}     % organize your images and other figures under media/ folder
\usepackage[authoryear]{natbib}
\usepackage{amsmath}

\usepackage{setspace}

\usepackage{xargs}                      % Use more than one optional parameter in a new commands
\usepackage[pdftex,dvipsnames]{xcolor}  % Coloured text etc.
\usepackage[colorinlistoftodos,prependcaption,textsize=small]{todonotes}

%\pgfplotsset{compat=1.18}
\newcommandx{\unsure}[2][1=]{\todo[linecolor=red,backgroundcolor=red!25,bordercolor=red,#1]{#2}}

\newcommandx{\change}[2][1=]{\todo[linecolor=blue,backgroundcolor=blue!25,bordercolor=blue,#1]{#2}}

\newcommandx{\info}[2][1=]{\todo[linecolor=OliveGreen,backgroundcolor=OliveGreen!25,bordercolor=OliveGreen,#1]{#2}}

\newcommandx{\improvement}[2][1=]{\todo[linecolor=Plum,backgroundcolor=Plum!25,bordercolor=Plum,#1]{#2}}

\newcommandx{\thiswillnotshow}[2][1=]{\todo[disable,#1]{#2}}

\newcommand{\bahar}[1]{\vspace{.1in}\todo[inline,color=green!50!white]{\textbf{Bahar:} #1}}

\newcommand{\question}[1]{\vspace{.1in}\todo[color=red!80!white]{\textbf{Question:} #1}}

\newcommand{\james}[1]{\vspace{.1in}\todo[inline,color=purple!20!white]{\textbf{James:} #1}}

\newcommand{\stephen}[1]{\vspace{.1in}\todo[inline,color=blue!20!white]{\textbf{Stephen:} #1}}



%Header
\pagestyle{fancy}
\thispagestyle{empty}
\rhead{ \textit{ }} 

% Update your Headers here
\fancyhead[LO]{Distribution-Free TSP Estimator for Arbitrary Dimension}
% \fancyhead[RE]{Firstauthor and Secondauthor} % Firstauthor et al. if more than 2 - must use \documentclass[twoside]{article}
\usepackage{xcolor}



  
%% Title
\title{Extending Distribution-Free Estimators for the Traveling Salesman Problem to Arbitrary Dimensions 
%Distribution-Free Estimators for the Traveling Salesman Problem on Arbitrary Dimension
%%%% Cite as
%%%% Update your official citation here when published 
%\thanks{\textit{\underline{Citation}}: 
%\textbf{Authors. Title. Pages.... DOI:000000/11111.}} 
}

\author{
  Stephen Feldman, James P. Bailey, Bahar \c{C}avdar \\
  Rensselaer Polytechnic Institute, Troy, NY\\
  \texttt{\{feldms5, bailej6, cavdab2\}@rpi.edu} \\
  %% examples of more authors
  % \And
  %Author3 \\
  %Affiliation \\
  %Univ \\
  %City\\
  %\texttt{email@email} \\
  %% \AND
  %% Coauthor \\
  %% Affiliation \\
  %% Address \\
  %% \texttt{email} \\
  %% \And
  %% Coauthor \\
  %% Affiliation \\
  %% Address \\
  %% \texttt{email} \\
  %% \And
  %% Coauthor \\
  %% Affiliation \\
  %% Address \\
  %% \texttt{email} \\
}


\begin{document}
\maketitle


\begin{abstract}
The Traveling Salesman Problem (TSP) remains a cornerstone of combinatorial optimization.
In many cases, we only need an estimation of the optimal tour length rather than solving for the route. 
Such cases can arise in strategic decisions, e.g., facility location, accounting for the cost of the future TSP routes. 
Hence, the ability to compute accurate TSP tour-length estimates across various scenarios in a short time is important.
Research on TSP tour-length estimation models has advanced in recent years, introducing new generation estimation models.
However, there are still two main limitations. 
\textcolor{red}{What are the current limitations?}.
To address these challenges, we present a new TSP tour-length estimation model based on using machine learning to predict the relationship between the length of the Minimum Spanning Tree (MST)  and the TSP tour-length. 
\textcolor{red}{approach to develop the model + contributions}.
%yet the computational cost of obtaining exact or near-optimal solutions often necessitates the use of tour length estimators. While theoretical bounds like the Beardwood-Halton-Hammersley (BHH) theorem provide asymptotic estimates, they rely on geometric assumptions—such as uniform point distribution and low dimensionality—that degrade rapidly in real-world, anisotropic, or high-dimensional environments. 
To address these limitations, we propose the Geometrically Assisted Regression Tree (GART), a machine learning framework that decomposes tour-length estimation into two components: the deterministic computation of the Minimum Spanning Tree (MST) and the statistical prediction of the MST-to-TSP scaling factor ($\alpha$).
By avoiding coordinate-dependent metrics such as the convex hull, GART maintains $O(n^2)$ computational complexity and eliminates the ``curse of dimensionality" inherent in geometric estimators. We evaluate GART against alternative estimators on complex, anisotropic hyperspaces. Experimental results demonstrate that GART achieves superior accuracy (lower Mean Absolute Percentage Error) and robustness across diverse distributions and dimensions %($D \in [2, 100]$) 
while offering a strict asymptotic speedup over constructive heuristics such as the Christofides algorithm. The model generalizes well outside of the training data, scaling to higher dimensions and additional customers. The proposed framework effectively bridges the gap between theoretical bounds and operational reality, enabling high-fidelity, millisecond-scale tour estimation for large-scale routing systems.

%GART utilizes a Gradient Boosted Decision Tree (LightGBM) regressor trained on a novel, dimension-agnostic feature set derived  from the graph topology of the MST, including metrics for branching factors, spectral diameter, and edge weight distribution moments. By avoiding coordinate-dependent metrics like the Convex Hull, GART maintains $O(N^2)$ computational complexity—dominated solely by the distance matrix calculation—and eliminates the "curse of dimensionality" inherent in geometric estimators. We evaluate GART against alternative estimators, including BHH, Chien's Probabilistic Model, and Space-Filling Curve heuristics, on complex, anisotropic hyperspaces. Experimental results demonstrate that GART provides superior accuracy (lower Mean Absolute Percentage Error) and robustness across diverse distributions ($D \in [2, 100]$) while offering a strict asymptotic speedup over constructive heuristics like the Christofides algorithm. The model generalizes well outside of the training data, scaling to higher dimensions and additional customers. The proposed framework effectively bridges the gap between theoretical bounds and operational reality, enabling high-fidelity, millisecond-scale tour estimation for large-scale routing systems.
\end{abstract}


% keywords can be removed
\keywords{Traveling Salesman Problem\and Tour length estimation \and Distribution-free \and Multidimensional networks}



\onehalfspacing

\section{Introduction}

The Traveling Salesman Problem (TSP) is a fundamental model in combinatorial optimization, widely used for logistics operations planning, network design, and resource allocation \citep{chien1992operational, concorde2006}. Because determining exact solutions for large-scale instances is NP-hard, the ability to compute rapid optimal tour-length estimates is essential for real-time decision-making \citep{heldkarp1962, wscube2024}.
In addition, these estimates can be used in other problems in which the TSP is formulated as a subproblem to improve computational efficiency (e.g., mobile facility location problem \citep{halper2011mobile}).



\subsection{Background and Motivation}

In the context of modern commerce, routing optimization radically improves supply chain efficiency. Last-mile delivery alone accounts for approximately 53\% of total shipping costs \citep{finmile2025}. Global logistics leaders have demonstrated that even marginal improvements in routing yield significant returns; for instance, the implementation of the UPS ORION system saves 100 million miles and \$400 million annually \citep{holland2017ups}. Furthermore, advanced route optimization software, such as Esri's ArcGIS Network Analyst, enables organizations to reduce operational costs by 10--30\% through optimized fleet management \citep{esri2023network}. Similarly, total route counts for major retailers can be reduced by 30--40\% through strategic route planning \citep{finmile2025}.

Beyond traditional logistics, fast TSP length approximation is a critical component across diverse applications in which TSP-like structures arise as subproblems. One of these application areas is bioinformatics. The ordering of markers in radiation hybrid mapping and the assembly of DNA fragments in genome sequencing are modeled as TSPs, where rapid cost estimation enables evaluation of assembly quality without full reconstruction \citep{agarwala2000fast, manfrin2006parallel}.
Moreover, in data compression, methods such as the Burrows-Wheeler Transform often require reordering data blocks to minimize transition costs, a problem analogous to finding a shortest path through a high-dimensional state space \citep{lloyd1982}. 
Additionally, in sensor network design, estimating the minimum data collection path length is crucial for optimizing the lifespan of battery-constrained mobile sinks in wireless sensor networks \citep{wang2008}.

In these domains, the problem is rarely confined to 2D Euclidean space. For example, genome assembly operates in abstract metric spaces defined by overlap similarity, while data compression and sensor networking often involve high-dimensional feature vectors. In such fields, metaheuristics often partition massive global problems into smaller clusters, each solved as an independent TSP to improve convergence speed \citep{frontiers2021, aaai2023}. Fast approximations of these sub-TSP instances allow algorithms to rapidly prune the search space and focus computational resources on promising regions.

\textcolor{red}{In the introduction, we mention the dimension gap as the only gap we are filling in. Is this how we want to position the paper? Addresses gap created by relying only on geometric properties introduced by.... E.g., we are more heavily relying on connections to MST which more effectively captures outliers and disjoint clusters...... Shoudl also mention that we outperform exiting benchmarks}





\textcolor{red}{Do we want the subsections in the introduction?}

\subsection{Limitations of the Current Models}
While the need for a TSP estimation model for general cases is clear, a persistent gap remains in the literature regarding high-dimensional or non-Euclidean instances. Traditional estimators, such as the BHH theorem \citep{bhh1959} and the probabilistic model by \cite{chien1992operational}, were developed primarily for 2D Euclidean space. 
Moreover, these models commonly rely on the calculation of the convex hull \citep{preparata1977} to estimate the ``volume" or ``area" of a point set to predict tour length. However, computing the convex hull becomes computationally prohibitive as the dimensionality increases \citep{menger1928}. 
\textcolor{orange}{Furthermore, it is possible to estimate the TSP tour length in problems from non-Euclidean metric spaces that nonetheless satisfy the triangle inequality, by converting them to a high-dimensional Euclidean metric space \citep{menger1928}.} \stephen{Stephen will move this sentence. }


\subsection{Contribution: The GART Model}
Our work contributes to the literature on distribution-free estimators, models that do not require prior knowledge of the node distribution, first introduced by \citet{cavdar2015distribution}. We advance this methodology by removing dependencies on features that scale poorly with dimensionality, such as convex hulls and Principal Component Analysis (PCA) \citep{jolliffe2002}. 
We introduce new features and develop a model to approximate the TSP tour length in $n$-dimensional spaces.  

The primary contribution of this work is the Geometrically Assisted Regression Tree (GART) model, a machine learning framework built utilizing Light  Gradient Boosted Decision Tree (LightGBM) \citep{ke2017lightgbm}. Unlike previous models that predict the tour length based on area, GART uses the relationship between the Minimum Spanning Tree (MST) and the optimal TSP tour. 
%predicts the $\alpha$ ratio: the multiplicative relationship between the Minimum Spanning Tree (MST) \citep{prim1957} and the optimal TSP tour length. 
By utilizing MST-centric features, such as the leaf ratio, dominance ratio, and normalized diameter, GART captures the underlying topological structure of a TSP instance. This enables the model to remain dimension-agnostic and avoid the limitations of geometric feature selection used in prior machine learning approaches.

\subsection{Paper Outline}

The remainder of this paper is structured as follows. Section \ref{sec:theory} introduces the theoretical foundations of the multiplicative relationship between the MST and TSP, which is analytically tightened through Christofides' algorithm \citep{christofides2022worst}. 
Section \ref{sec:methodology} describes the construction of the GART model, including feature selection, construction of our $n$-dimensional training dataset utilizing Concorde TSP solver \citep{concorde2006} and LKH-3 algorithm \citep{helsgaun2017}. Section \ref{sec:model} presents the trained model with hyperparameter tuning via Optuna \citep{akiba2019optuna}, and analysis of model complexity. In Section \ref{sec:evaluation}, we present a survey of existing tour length estimators we benchmark GART against and our comparison results. 
%a construction of an adversarial 2-dimensional dataset, and the results.
Section \ref{sec:dsicussion} further analyzes the benchmarking results and compares GART to other deep learning approaches. Finally, we conclude our paper with Section \ref{sec:conclusion}, including future research directions. 

\section{Theoretical Foundations} \label{sec:theory}

To formally situate the estimation problem, we first define the Traveling Salesman Problem (TSP) using its classic integer linear programming (ILP) formulation. Let $G = (V, E)$ be a complete undirected graph where $V = \{1, \dots, n\}$ is the set of vertices and $E = \{(i, j) : i, j \in V, i < j\}$ is the set of edges. Associated with each edge $(i, j)$ is a non-negative cost $c_{ij}$, typically representing the Euclidean distance between points in $\mathbb{R}^d$. The objective is to find a tour of minimum total cost that visits every vertex exactly once. Following the standard Dantzig-Fulkerson-Johnson formulation found in \citet{concorde2006}, we use binary decision variables $x_{ij}$, where $x_{ij} = 1$ if edge $(i, j)$ is included in the tour, and $0$ otherwise. The problem formulation is as follows:
\begin{equation} \label{eq:tsp_obj}
    \min \sum_{i=1}^{n} \sum_{j=i+1}^{n} c_{ij} x_{ij}
\end{equation}

Subject to the degree constraints:
\begin{equation} \label{eq:degree}
    \sum_{j < i} x_{ji} + \sum_{j > i} x_{ij} = 2, \quad \forall i \in V
\end{equation}

And the subtour elimination constraints:
\begin{equation} \label{eq:subtour}
    \sum_{i \in S} \sum_{j \in S, j > i} x_{ij} \le |S| - 1, \quad \forall S \subset V, 2 \le |S| \le n-2
\end{equation}

%The combinatorial explosion of subtour elimination constraints renders exact solution methods, such as branch-and-cut, computationally intractable for large-scale real-time applications, as the problem is NP-hard \citep{heldkarp1962}.

The TSP is NP-hard \citep{heldkarp1962}.
In contrast to the intractable TSP, the Minimum Spanning Tree (MST) problem, which seeks a connected, acyclic subgraph minimizing the total edge weight, can be solved optimally in polynomial time.
%MST seeks a connected, acyclic subgraph minimizing the total edge weight. 
%For the same graph $G$, the MST can be solved optimally in polynomial time. 
Algorithms such as those proposed by \citet{prim1957} construct the MST with a time complexity of $O(E + V \log V)$, or $O(N^2)$ for complete graphs. 
Crucially, the computational cost of the MST scales linearly with the dimensionality $d$ of the edge weights $c_{ij}$, making it a computationally stable baseline for high-dimensional inference.

The relationship between these two problems provides a theoretical bound on the optimal TSP tour length, denoted $L_{\text{TSP}}$. Since any valid TSP tour is a connected graph spanning all vertices, removing one edge yields a spanning tree. Thus, the weight of the MST, $L_{\text{MST}}$, is a strict lower bound. Conversely, a valid tour can be constructed by traversing the MST edges twice (a depth-first walk) and applying shortcuts via the triangle inequality, establishing the classic double-tree bound:
\begin{equation} \label{eq:bounds}
    L_{\text{MST}} \le L_{\text{TSP}} \le 2 \cdot L_{\text{MST}}.
\end{equation}

\citet{christofides2022worst} improved this upper bound by augmenting the MST with a Minimum Weight Perfect Matching (MWPM) on the set of vertices with odd degrees. This yields a tour length guaranteed to be at most $1.5 \cdot L_{\text{TSP}}$ in metric spaces. However, calculating the MWPM requires $O(N^3)$ time, which is significantly more expensive than the $O(N^2)$ MST construction. For large-scale auxiliary subproblems where computation speed is critical, the cubic complexity of the Christofides refinement is often prohibitive, necessitating a faster method to approximate the tightening coefficient.

Leveraging the relationship between the optimal TSP and MST values, instead of explicitly constructing the matching, we treat the tightening coefficient as a statistical parameter $\alpha = L_{\text{TSP}} / L_{\text{MST}}$. 
While the worst-case bound for $\alpha$ is $1.5$, its expected behavior is heavily dependent on the dimensionality $d$. \citet{percus1996} utilized statistical mechanics to demonstrate that in the limit of high dimensions ($d \to \infty$), the lengths of the MST and the TSP converge, driving the ratio $\alpha \to 1$. 
This implies that the ``cost" of transforming a spanning tree into a TSP tour diminishes as the spatial freedom increases. 
Consequently, a static estimator calibrated for 2D (where $\alpha \approx 1.075$) will systematically overestimate costs in higher dimensions. 
Our proposed model addresses this by learning to predict $\alpha$ dynamically based on the topological characteristics of the node set.


\section{Methodology: The Geometrically Assisted Regression Tree} \label{sec:methodology}

We propose a new method to estimate the TSP tour length for general dimensionality levels, using the relationship between the MST. Our method, which we call Geometrically Assisted Regression Tree (GART), integrates the regression tree approach with \textcolor{red}{\{specific to our method\}} to overcome the dimensional limitations of classical geometric estimators.
%We propose the Geometrically Assisted Regression Tree (GART), a framework designed to bypass the dimensional limitations of classical geometric estimators. 
Unlike traditional approaches that map coordinate-dependent metrics directly to tour length, GART decouples the estimation into two distinct phases: the calculation of the Minimum Spanning Tree (MST) and the statistical prediction of the scaling coefficient $\alpha$, denoted by $\hat{\alpha}$. The final estimator is defined as $\hat{L}_{\text{TSP}} = \hat{\alpha} \cdot L_{\text{MST}}$, where $\hat{\alpha}$ is predicted by a gradient-boosted decision tree regressor based on the topological characteristics of the instance. The following section introduces the features we use to capture these characteristics.
%topological signature of the instance.

\subsection{Feature Selection} \label{subsec:features}

To predict the scaling factor $\alpha$ accurately across varying dimensions, we extract a comprehensive set of 28 features from each TSP instance. \bahar{we should explain the number of features since there are quite many.}
These features are designed to be dimension-agnostic, relying primarily on the graph topology of the Minimum Spanning Tree (MST) and the relative distribution of points, rather than absolute coordinate positions. \bahar{are the current features in the literature based on absolute positions? I think they are also on relative distribution?}
\bahar{is dimensionality not a part of the topology?}
This approach ensures that the computational complexity of feature extraction remains dominated by the $O(N^2)$ construction of the MST, avoiding the exponential costs associated with high-dimensional geometric partitioning. \stephen{this approach alone does not ensure it is faster than $o(n^2)$. do we ensure it by choosing the features that are faster than $o(n^2)$?}
We categorize the features into three groups: (1) geometric and centroid descriptors, (2) MST edge distribution statistics, and (3)topological structure and clustering proxy, as described below in detail. \bahar{we may need some explanation for the categories here.}



\subsubsection{Geometric and Centroid Descriptors}

The first category of features quantifies the macroscopic spatial properties of the point cloud. \bahar{what is point cloud? do we mean the node set in $R^d$?}
\stephen{New convention: Convert all points/}
Geometric descriptors provide the \textcolor{red}{baseline scaling parameters required by asymptotic bounds}, such as the BHH theorem \citep{bhh1959}. \bahar{this sentence is not clear to me - what are we referring to?}
Centroid statistics, adapted from \citet{cavdar2015distribution}, measure the dispersion of nodes relative to the geometric center of the instance, enabling the model to distinguish between different node dispersions, e.g., uniform distributions and centralized Gaussian clusters.
In Table \ref{tab:geo_features}, we present the list and the description of these metrics. % details these metrics.

\bahar{The wording should be consistent. We defined the graph G in Section 2, composed of set of vertices and set of edges. We also say points (interchangeably for the edges, I think), but it is not introduced. We should decide which is best in this context.}
\begin{table}[!ht]
\centering
\caption{Geometric and Centroid Feature Specifications.}
\label{tab:geo_features}
\resizebox{\textwidth}{!}{%
\begin{tabular}{@{}lp{6.5cm}ll@{}}
\toprule
\textbf{Feature Name} & \textbf{Description \& Rationale} & \textbf{Complexity} & \textbf{Origin / Citation} \\ \midrule
\textbf{Dimension ($D$)} & The number of spatial coordinates defining each point. & $O(1)$ & Standard \\
\textbf{Number of Customers ($N$)} & The total count of cities in the instance. & $O(1)$ & Standard \\
\textbf{Aspect Ratio} & Ratio of the largest dimension range to the smallest. Detects subspace distortion. & $O(D)$ & Standard \\
\textbf{Bounding Hypervolume} & Product of coordinate ranges: $\prod (max_d - min_d)$. Approximates volume $V$ for BHH bounds. & $O(N \cdot D)$ & \citet{bhh1959} \\
\textbf{Node Density} & Ratio of $N$ to Bounding Hypervolume. Proxy for point intensity. & $O(1)$ & \citet{smithmiles2010} \\
\textbf{Centroid Dist (Mean, Std)} & Average and deviation of distances to the centroid. Measures central tendency. & $O(N \cdot D)$ & \citet{cavdar2015distribution} \\
\textbf{Centroid Dist Max} & Radius of the enclosing sphere. & $O(N \cdot D)$ & \citet{cavdar2015distribution} \\
\textbf{Centroid Dist IQR} & Interquartile range of centroid distances. Robust dispersion metric. & $O(N \cdot D)$ & \citet{cavdar2015distribution} \\ \bottomrule
\end{tabular}%
}
\end{table}


\james{Should be clear when these features are modified from other papers (in text and in table)... e.g., IQR is not in \cite{cavdar2015distribution}.}

\subsubsection{MST Edge Distribution Statistics}

The second category captures the local uniformity of the point set through the statistical moments of the MST edge weights. Unlike coordinate-based variance, edge weights are invariant to rotation and translation. High skewness or kurtosis in this distribution typically indicates a mixture of scales, such as dense clusters separated by voids, which significantly impacts the $\alpha$ ratio.

\begin{table}[!ht]
\centering
\caption{MST Edge Statistic Specifications}
\label{tab:edge_features}
\resizebox{\textwidth}{!}{%
\begin{tabular}{@{}lp{6.5cm}ll@{}}
\toprule
\textbf{Feature Name} & \textbf{Description \& Rationale} & \textbf{Complexity} & \textbf{Origin / Citation} \\ \midrule
\textbf{MST Edge Max} & Weight of the longest edge in the MST. & $O(N)$ & Standard \\
\textbf{MST Quantiles} & 10\%, 25\%, 50\%, 75\%, 90\% percentiles. Provides a non-parametric distribution profile. & $O(N \log N)$ & Standard Statistical \\
\textbf{MST Total Length} & Sum of all edge weights. The primary scaling factor; $L_{TSP} \ge L_{MST}$. & $O(N^2)$ & \citet{prim1957} \\
\textbf{MST Edge Mean} & Average edge length in the MST. & $O(N)$ & \citet{smithmiles2010} \\
\textbf{MST Edge Std} & Standard deviation of edge lengths in the MST. & $O(N)$ & \citet{smithmiles2010} \\
\textbf{MST Edge Skewness} & Asymmetry of edge distribution. High positive skew implies many short edges and few long bridges. & $O(N)$ & \citet{smithmiles2010} \\
\textbf{MST Edge Kurtosis} & Tailedness of distribution. Indicates the presence of extreme outliers. & $O(N)$ & \citet{smithmiles2010} \\  \bottomrule
\end{tabular}%
}
\end{table}

\bahar{is there a reason for the ordering in the table? if not, can we collect the proposed ones together?}
\stephen{The current order was grouping by similar measures - centroid distance goes together etc, but I have no objections to structuring proposed versus existing}
\subsubsection{Topological Structure and Clustering Proxies}

The final category contains features designed to detect higher-order structures—such as clusters, filaments, and hubs—without performing expensive clustering algorithms. We introduce specific ratios (Dominance, Gap, Normalized Diameter) that serve as $O(N)$ proxies for these complex geometric properties. For instance, the "Leaf Ratio" effectively distinguishes between low-dimensional linear manifolds (low leaf count) and high-dimensional isotropic noise (high leaf count).

\begin{table}[!ht]
\centering
\caption{Topological and Clustering Feature Specifications}
\label{tab:topo_features}
\resizebox{\textwidth}{!}{%
\begin{tabular}{@{}lp{6.5cm}ll@{}}
\toprule
\textbf{Feature Name} & \textbf{Description \& Rationale} & \textbf{Complexity} & \textbf{Origin / Citation} \\ \midrule
\textbf{MST Dominance Ratio} & Sum of largest $\sqrt{N}$ edges / Total Length. Captures weight of inter-cluster bridges. & $O(N)$ & \textbf{Proposed} \\
\textbf{MST Gap Ratio} & Max edge / Median edge. Measures magnitude of sparsity jumps. & $O(1)$ & \textbf{Proposed} \\
\textbf{MST Leaf Ratio} & Fraction of nodes with degree $k=1$. Topology proxy for dimensionality. & $O(N)$ & \citet{smithmiles2010} \\
\textbf{MST Degree (Mean, Std)} & Branching factor statistics. & $O(N)$ & \citet{smithmiles2010} \\
\textbf{MST Degree Max} & Maximum node degree. Identifies "hub" nodes common in High-D spaces. & $O(N)$ & \citet{smithmiles2010} \\
\textbf{MST Diameter} & Total weight of the longest path in the tree. & $O(N)$ & Graph Theory \\
\textbf{Normalized Diameter} & Diameter / Total Length. Scale-invariant shape proxy ($\approx 1.0$ for lines). & $O(1)$ & \textbf{Proposed} \\
\textbf{Large Edge Count} & Count of edges $> \mu + \sigma$. Secondary cluster indicator. & $O(N)$ & Statistical Process Control \\ \bottomrule
\end{tabular}%
}
\end{table}



\subsection{Dataset Generation} \label{subsec:dataset}
\bahar{Do we need this section here? Can we move it to the Experimental Evaluation Section?}
\stephen{Yes we can move it.}
To train a general model across hyperspaces, we constructed a large-scale synthetic dataset with rich topological diversity. 
Unlike traditional benchmarks that often use a single distribution for node coordinates (e.g., placing points uniformly in a unit square), our instance generation process employs a ``Mix-and-Match" strategy. \textcolor{red}{
For each instance of dimension $D$, we generate a distribution signature string $S$ of length $D$, where $S_j$ corresponds to the probability distribution code for the $j$-th dimension. 
This methodology allows for the construction of complex anisotropic spaces where the topology varies by axis—for example, a 3-dimensional instance might follow a ``Unk'' pattern, where the x-axis is \textbf{U}niform, the y-axis is \textbf{N}ormal, and the z-axis is \textbf{K}orrelated.} \bahar{wording. I do not follow this paragraph. What is an unk pattern?}
\stephen{this belongs in the appendix is the main issue I think. This explains how to do in depth analysis per distribution on my data and a bit about how the distributions work. Might be worth dropping the conventions entirely. Its clear from reading the code.}


\textcolor{red}{
To ensure the model does not overfit to a limited set of random seeds, we employ a ``Genotypic Seeding'' strategy. The random seed for each instance is derived by hashing its unique distribution signature string. To maximize the entropy of these seeds while using only four primary distribution classes, we utilize a 15-character alphabet where multiple codes serve as functional aliases (see Table \ref{tab:distributions}). For instance, both `n''and `i' map to the Normal distribution function.
However, the signature ``rn'' will hash to a completely different random seed than ``ri'', resulting in distinct point realizations despite sharing the same statistical properties. This technique allows us to generate a combinatorially vast number of unique problem instances without introducing bias toward any specific geometric type.}



\bahar{Do we need this level of technical detail for instance generation? We can state that the instances are generated to avoid any bias toward any distribution type. I think this level of detail can either be omitted or moved to the appendix.}
\stephen{ref above comment, i agree, will move}

The training dataset consists of instances with node counts $N \in [20, 1000]$ and dimensions $D \in [2, 50]$. We deliberately restrict the training data to a maximum of 50 dimensions to test the model's ability to extrapolate to higher dimensions (up to $D=100$) during the evaluation phase.

\begin{table}[!ht]
\centering
\caption{Probability Distributions and Alias Codes}
\label{tab:distributions}
\resizebox{0.8\textwidth}{!}{%
\begin{tabular}{@{}llp{8cm}@{}}
\toprule
\textbf{Distribution Class} & \textbf{Alias Codes} & \textbf{Mathematical Definition ($x \in [0, G]$)} \\ \midrule
\textbf{Uniform} & r, p, t, g, l, s, b, x, e, a, h & $x \sim \mathcal{U}(0, G)$. Represents standard unclustered noise. Aliases ensure high seed entropy for uniform spaces. \\
\textbf{Normal} & n, i & $x \sim \mathcal{N}(G/2, G/6)$, clipped to $[0, G]$. Represents centralized blobs or single-cluster instances. \\
\textbf{Clustered} & c, o & $x \sim \mathcal{N}(\mu_k, 0.05G)$ around $K \approx \sqrt{N}/4$ centers. Represents distinct neighborhoods. \\
\textbf{Correlated} & k & $x_i = x_{i-1} + \mathcal{N}(0, 0.1G)$. The coordinate depends on the previous dimension, creating diagonal or linear manifolds. \\ \bottomrule
\end{tabular}%
}
\end{table}

For each generated instance, the optimal tour length was computed using state-of-the-art solvers. 
Each problem instance was solved using Concorde \citep{concorde2006} and LKH-3 \citep{helsgaun2017}, and the shortest length was stored as the optimal solution. \textcolor{blue}{To address the limitations of the standard TSPLIB format, which mandates integer edge weights, we implemented a dynamic coordinate scaling mechanism to preserve floating-point precision during the solver phase. 
Standard rounding of small coordinate differences—frequently observed in high-density clusters or normalized spaces—can erode fine-grained structural details, effectively collapsing distinct points into zero-cost edges. 
To prevent this data erosion, our pipeline automatically calculates a scaling factor $\sigma$ based on the instance grid size $G$, normalizing the effective coordinate space to a target range of $[0, 10^4]$. Specifically, Euclidean distances are computed as $d_{ij} = \lfloor ||x_i - x_j||_2 \cdot \sigma + 0.5 \rfloor$, with a strict enforcement of $d_{ij} \ge 1$ to prohibit zero-weight edges between distinct nodes. The solvers operate on these scaled integer matrices, and the resulting optimal tour lengths are subsequently divided by $\sigma$ to recover the precise floating-point cost.}
\bahar{Is this to say that we multiply the float coordinates by a large number to achieve integer values? If so, this is unnecessary, and I suggest removing this.}
\stephen{The integer scaling is required for the data to be parsed as input into the solvers, esp Concorde. This is absolutely a requirement. I agree that it is again a technical detail that is not required.}

\section{Estimation Model Development and Complexity Analysis} 
\label{sec:model}

Having established the dimension-agnostic feature set, we proceed to the design of the predictive model itself. The Geometrically Assisted Regression Tree (GART) is implemented using the LightGBM gradient-boosting framework \citep{ke2017lightgbm}.
%While Deep Learning approaches—specifically Graph Neural Networks (GNNs)—have gained traction for constructing TSP tours, they are often computationally expensive to train and difficult to interpret. 
The reason for 
For the distinct task of scalar regression on tabular topological features, Gradient Boosted Decision Trees (GBDTs) offer superior sample efficiency and faster inference speeds.

\subsection{Architecture Selection and Training}
\bahar{title is not clear (algorithm selection). a better wording?}
\stephen{typo, should say architecture, as in model type}
We selected LightGBM over standard Random Forests or traditional GBDT implementations due to its unique "leaf-wise" tree growth strategy. Unlike level-wise algorithms that grow trees depth-first to maintain balance, LightGBM splits the leaf with the maximum delta loss, regardless of depth. This approach is mathematically advantageous for approximating the $\alpha$ ratio because the relationship between graph topology and tour length often exhibits sharp, non-smooth phase transitions—for example, the sudden drop in $\alpha$ when a random point cloud begins to form distinct clusters. Leaf-wise growth allows the model to allocate more complexity to these critical transition regions without wasting splits on simpler, linear parts of the manifold.

The model was trained to minimize the Root Mean Squared Error (RMSE) between the predicted ratio $\hat{\alpha}$ and the ground truth $\alpha$. To optimize the model's generalization capability, we employed the Optuna framework \citep{akiba2019optuna} to perform a Bayesian search of the hyperparameter space. The optimization utilized the Tree-structured Parzen Estimator (TPE) algorithm over 100 trials to tune the learning rate, tree structure, and regularization terms. The final model configuration, detailed in Table \ref{tab:hyperparams}, resulted in a robust ensemble of 4,431 trees. While the hyperparameter search allowed for up to 300 leaves per tree to capture complex interactions, the final trained model exhibited a maximum leaf count of 258, indicating that the constraint was sufficient to prevent unbridled overfitting while allowing necessary depth.


\begin{table}[!ht]
\centering
\caption{Final Model Hyperparameters and Ensemble Statistics}
\label{tab:hyperparams}
\begin{tabular}{llp{7cm}}
\toprule
\textbf{Parameter} & \textbf{Value} & \textbf{Description} \\ \midrule
\textit{Ensemble Size} & $4,431$ & Total number of decision trees in the final model. \\
\textit{Max Leaves (Actual)} & $258$ & Maximum leaves observed in any single tree. \\
Max Leaves (Constraint) & $300$ & Upper bound on leaves during optimization. \\
Learning Rate & $0.0180$ & Step size shrinkage used to prevent overfitting. \\
Feature Fraction & $0.416$ & Randomly select 41.6\% of features on each iteration. \\
Bagging Fraction & $0.561$ & Randomly select 56.1\% of data for each tree. \\
Bagging Freq & $3$ & Perform bagging every 3 iterations. \\
Min Child Samples & $45$ & Minimum data needed in a leaf. \\
Lambda L1 & $0.417$ & L1 regularization term on weights (lasso). \\
Lambda L2 & $4.58 \times 10^{-6}$ & L2 regularization term on weights (ridge). \\ \bottomrule
\end{tabular}
\end{table}

\subsection{Computational Complexity Analysis}

A primary motivation for GART is to provide high-quality estimates faster than constructive heuristics. To validate this, we analyze the worst-case time complexity of the GART pipeline compared to the Christofides algorithm, which provides the tightest analytical bound ($1.5 \cdot MST$). The total runtime of GART is the sum of the feature extraction time and the inference time.

The computational pipeline begins with the construction of the Minimum Spanning Tree (MST). For a dense complete graph with $N$ vertices, calculating the distance matrix requires $O(N^2 \cdot D)$ operations. Constructing the MST using Prim's algorithm takes $O(N^2)$ time. Once the MST is established, extracting topological features such as degrees, diameter, and leaf counts requires a linear traversal of the tree edges, taking $O(N)$. Calculating edge quantiles requires sorting the $N-1$ edges, which adds a term of $O(N \log N)$. Therefore, the feature extraction overhead is dominated entirely by the initial distance matrix calculation.

The inference complexity of the gradient boosted model is determined by the total number of split operations across the ensemble. Given our final model structure of $K=4,431$ trees with a maximum of $L=258$ leaves per tree, the inference requires at most $K \times L \approx 1.14 \times 10^6$ comparisons. Crucially, this value is a model constant independent of the problem input size $N$. Consequently, the inference time is $O(1)$ relative to $N$, and the total asymptotic complexity of the GART pipeline remains $O(N^2 \cdot D)$.

In contrast, the Christofides heuristic requires finding a Minimum Weight Perfect Matching (MWPM) on the odd-degree vertices. While the theoretical worst-case for exact MWPM algorithms is $O(N^3)$ \citep{cook1999}, recent approximation algorithms by \citet{duan2014} and \citet{drake2003} can achieve near-linear time complexity for the matching step.However, while near-linear time approximations for the matching step exist \citep{duan2014}, utilizing them fundamentally alters the theoretical guarantees of the heuristic. The classic 1.5-approximation ratio of the Christofides algorithm relies strictly on the exact minimality of the perfect matching; substituting an approximate matching algorithm inevitably relaxes this bound, often degrading the worst-case ratio to $1.5 + \epsilon$ or worse depending on the greedy heuristics employed \citep{drake2003}. Furthermore, asymptotic analysis often obscures the significant constant factors inherent in matching algorithms. Implementing sophisticated blossom-shrinking protocols or randomized matching data structures involves implementation overheads that far exceed the simple, highly parallelizable algebraic operations required by GART. 

Most critically, constructive heuristics like Christofides provide a worst-case upper bound, which is often overly pessimistic for real-world decision-making. In logistics contexts, such as dynamic pricing or time-window feasibility checks, relying on a loose upper bound (e.g., predicting $1.5 \times \text{MST}$ when the true cost is typically $1.08 \times \text{MST}$) leads to inefficient resource allocation and uncompetitive service guarantees. GART improves upon this by providing a statistically expected estimate rather than a theoretical bound, aligning the predicted cost more closely with the operational reality of the problem instance.

\section{Experimental Evaluation} \label{sec:evaluation}

To validate the efficacy of the Geometrically Assisted Regression Tree (GART), we conducted a comprehensive benchmarking study against a suite of established theoretical and constructive estimators. The evaluation methodology was designed to assess predictive fidelity, topological robustness, and computational efficiency across a diverse spectrum of problem instances. Unlike prior studies that often rely on uniform distributions, our testing framework incorporates pathological cases—including anisotropic clusters and linear filaments—to rigorously stress-test the asymptotic assumptions of geometric baselines.

\subsection{Benchmark Models}

We compared GART against five baseline estimators representing the evolution of TSP approximation theory. To ensure reproducibility, we present the exact mathematical formulations used in our Python implementation, with constants derived directly from the source code.

\subsubsection{Beardwood-Halton-Hammersley (BHH) Theorem}
The foundational asymptotic bound for TSP assumes the tour length scales with the hypervolume of the point set \citep{bhh1959}. Our implementation estimates the tour length $L$ as:
\begin{equation}
    \hat{L}_{\text{BHH}} = \beta_d \cdot N^{\frac{d-1}{d}} \cdot V^{\frac{1}{d}}
\end{equation}
where $N$ is the number of nodes, $V$ is the volume of the convex hull (or bounding box in high dimensions), and $\beta_d$ is a dimension-dependent constant. For $d=2$, we utilize $\beta_2 \approx 0.7124$.

\subsubsection{Distribution Free Estimator }
To handle anisotropic distributions, \citet{cavdar2015distribution} proposed a distribution-free estimator that replaces simple volume metrics with statistical moments of coordinate deviations. This estimator is restricted to 2D with a forula:
\begin{equation}
    \hat{L}_{\text{CS}} = \alpha_1 \sqrt{N \sigma'_x \sigma'_y} + \alpha_2 \sqrt{\frac{N A \sigma_x \sigma_y}{\bar{c}_x \bar{c}_y}}
\end{equation}
where $A$ is the area of the convex hull, $\sigma_k$ is the standard deviation of coordinates in dimension $k$, $\bar{c}_k$ is the mean absolute deviation from the centroid, and $\sigma'_k$ is the standard deviation of these absolute deviations. The regression coefficients are fitted by \citet{cavdar2015distribution} as $\alpha_1 = 2.791$ and $\alpha_2 = 0.2669$.

\subsubsection{Hilbert Space-Filling Curve}
This constructive heuristic maps $N$-dimensional points onto a 1-dimensional interval $[0, 1]$ using a fractal Hilbert curve \citep{bartholdi1982heuristic}. The points are sorted by their Hilbert index $h(x)$, and the tour is constructed by visiting nodes in this sorted order:
\begin{equation}
    \hat{L}_{\text{Hilbert}} = \sum_{i=1}^{N-1} ||x_{(i)} - x_{(i+1)}||_2 + ||x_{(N)} - x_{(1)}||_2
\end{equation}
This approach preserves locality without constructing a full distance matrix, offering $O(N \log N)$ complexity.

\subsubsection{MST Ratio}
The control baseline scales the weight of the Minimum Spanning Tree by a dimension-dependent factor $\rho(d)$ derived from the Prim's algorithm lower bound \citep{prim1957}:
\begin{equation}
    \hat{L}_{\text{MST-Ratio}} = \rho(d) \cdot L_{\text{MST}}
\end{equation}
For 2D instances, we utilize $\rho(2) = 1.075$.

%Section on the 2d dataset added back in.
\subsection{Adversarial 2D Dataset Construction} \label{subsec:dataset_construction}

To rigorously evaluate the GART framework, we generated a bespoke benchmark suite comprising 2,580 distinct instances. While standard datasets often default to uniform distributions where asymptotic bounds hold trivially, this suite is designed to expose specific failure modes in geometric estimation by varying the point set topology $N \in [5, 1000]$ and grid scaling $G \in [1000, 10000]$.

The dataset is partitioned into five distinct classes, each targeting a specific assumption inherent in classical BHH-type estimators. Table \ref{tab:dataset_counts} provides the complete breakdown of instance counts and the adversarial characteristics of each class.

\begin{table}[!ht]
\centering
\caption{2D Benchmark Dataset Composition ($N_{total} = 2,580$)}
\label{tab:dataset_counts}
\resizebox{\textwidth}{!}{%
\begin{tabular}{@{}llcp{6cm}@{}}
\toprule
\textbf{Distribution Class} & \textbf{Generators} & \textbf{Count} & \textbf{Adversarial Characteristic} \\ \midrule
\textbf{Isotropic Noise} & Random, Triangular, Truncated Exp. & 630 & \textit{Maximal Convergence.} Represents the ideal case for Area-based estimators ($A_{Hull} \approx A_{Box}$). \\ \midrule
\textbf{Biased / Skewed} & Squeezed, Uni-Tri, Tri-Squeezed, Correlated & 840 & \textit{Aspect Ratio Distortion.} Challenges estimators assuming equal variance ($\sigma_x \approx \sigma_y$). \\ \midrule
\textbf{Geometric Struct.} & Grid, Boundary (Hollow), X-Central & 630 & \textit{Regularity.} Tests detection of lattice structures and hollow voids that reduce entropy. \\ \midrule
\textbf{Clustered} & Variable Radius $k$-Means ($k \in [5, 20]$) & 60 & \textit{Density Variance.} Violates uniform density assumption ($N/A$). Introduces bimodal edge weights. \\ \midrule
\textbf{Line Noise} & Linear Manifold ($y=mx+b$) + $\epsilon$ & 210 & \textit{Area Stress Test} Forces $A_{BoundingBox}$ farther from $A_{Convexhull}$. Designed to challenge GART. \\ \bottomrule
\textbf{Total} & & \textbf{2,580} & \\ \bottomrule
\end{tabular}%
}
\end{table}

\subsubsection{The Geometric Enclosure Dilemma: Hull vs. Box}
A critical motivation for this varied dataset is to interrogate the reliability of the "area" metric itself. Classical estimators derive the expected nearest-neighbor distance—and thus the tour length—from the point density $\delta = \frac{N}{V}$, where $V$ is the volume of the enclosing region. However, defining this region presents a fundamental dilemma that is exacerbated in high-dimensional or anisotropic data: the choice between the Bounding Box ($V_{BB}$) and the Convex Hull ($V_{CH}$).

For isotropic distributions (e.g., Uniform), the convex hull approximates the bounding box ($V_{CH} \approx \frac{1}{2} V_{BB}$ in 2D), making either a valid proxy for density. However, for anisotropic distributions such as the \textit{Correlated} or \textit{Line Noise} classes included in our benchmark, this relationship diverges radically. As the points align along a diagonal manifold, the volume of the convex hull shrinks ($V_{CH} \to 0$) while the bounding box volume remains static and large. 

An estimator relying on $V_{BB}$ will perceive a "large space" with "low density," erroneously predicting a short tour length dominated by empty corners. Conversely, an estimator using $V_{CH}$ correctly identifies the high effective density along the manifold. While using the Convex Hull provides the tightest geometric bound, calculating it becomes computationally intractable as dimensions increase ($O(N^{\lfloor d/2 \rfloor})$).

This dataset explicitly tests GART's "geometry-free" hypothesis. Since GART discards coordinate information in favor of MST topology, it lacks access to the explicit Convex Hull volume. The \textbf{Isotropic Noise} subset (630 instances) serves as a control group to measure the "cost of blindness": does GART suffer a performance penalty on well-behaved instances where the Convex Hull is the mathematically optimal predictor? Conversely, the \textbf{Line Noise} and \textbf{Biased} subsets test whether GART's topological features (e.g., Normalized Diameter, Leaf Ratio) can implicitly recover the density information that Bounding Box estimators miss, without incurring the computational cost of the Convex Hull.

\subsubsection{The Line Noise Singularity}
The \textit{Line Noise} instances represent the limit case of the enclosure dilemma. As the width of the point distribution $\epsilon \to 0$, the area of the convex hull vanishes ($A \to 0$). Geometric formulas dependent on $\sqrt{A}$, such as the BHH theorem, mathematically collapse to a predicted cost of zero, despite the tour length remaining strictly positive (approximating $2 \times \text{Length}_{line}$). By including 210 such instances, we evaluate whether GART's regression tree can detect this dimensional collapse via graph-theoretic signals—specifically the convergence of the MST Leaf Ratio to 2.0—and switch its internal logic from area-scaling to linear summation, a capability that strict geometric formulas lack.

%Benchmark results
\subsection{Performance Results and Analysis}

Table \ref{tab:unified_results} presents the unified performance metrics. GART achieves the highest accuracy on the 2D benchmark ($R^2=0.998$) and demonstrates superior generalization on the N-Dimensional set.


\james{Is N the dimension or number of nodes? Check everywhere}

\james{Can we see something like the number of nodes vs ratios? Currently, we are lumping everything together, which loses nuance. }

\bahar{CIs? Do you want paired CIs to control for variability? Maybe ... E.g., to pair: $GART_i$=error for Gart for instance i, $HILBERT$= error for instance i.  CI on $GART_i/HILBERT_i$. Unclear if needed. }

\begin{table}[!ht]
\centering
\caption{Unified Performance Matrix: 2D Benchmarks vs. N-Dimensional Extrapolation}
\label{tab:unified_results}
\resizebox{0.7\textwidth}{!}{%
\begin{tabular}{@{}lccccc@{}}
\toprule
\textbf{Metric} & \textbf{GART} & \textbf{MST Ratio} & \textbf{Hilbert} & \textbf{Cavdar} & \textbf{BHH} \\ \midrule
\multicolumn{6}{l}{\textit{\textbf{2D Performance ($N \le 1000$)}}} \\ 
$R^2$ Score & \textbf{0.998} & 0.992 & 0.801 & 0.909 & 0.943 \\
MAPE (\%) & \textbf{3.23} & 12.45 & 31.01 & 19.82 & 23.82  \\
Bias (\%) & \textbf{0.21} & -11.76 & 31.01 & 11.42 & -13.66  \\
Stability ($\sigma$) & \textbf{4.94} & 11.06 & 14.24 & 42.04 & 32.05 \\ \midrule
\multicolumn{6}{c}{\textbf{Avg. 2D Optimal Solver Time: 9.53 s}} \\ \midrule
\multicolumn{6}{l}{\textit{\textbf{2D Efficiency ($N \le 1000$)}}} \\ 
Feature Time (s) & 0.132 & 0.041 & 0.015 & 0.007 & 0.006  \\
Inference Time (s) & 0.038 & 0.041 & 0.015 & 0.007 & 0.006  \\
Total Time (s) & 0.170 & 0.081 & 0.030 & 0.013 & 0.011  \\
Speedup (\%) & 48.74 & 90.20 & 97.71 & 95.47 & 94.76 \\ \midrule \midrule
\multicolumn{6}{l}{\textit{\textbf{N-Dimensional Performance ($D \in [3, 100]$)}}} \\ 
$R^2$ Score & \textbf{0.9987} & 0.9992 & 0.9621 & \multicolumn{2}{c}{\textit{-- Failed --}} \\
MAPE (\%) & \textbf{0.81} & 8.56 & 21.28 & \multicolumn{2}{c}{\textit{-- Failed --}} \\
Bias (\%) & \textbf{0.31} & -8.56 & 21.28 & \multicolumn{2}{c}{\textit{-- Failed --}} \\
Stability ($\sigma$) & \textbf{1.17} & 6.89 & 14.56 & \multicolumn{2}{c}{\textit{-- Failed --}} \\ \midrule
\multicolumn{6}{c}{\textbf{Avg. N-Dim Optimal Solver Time: 1.72 s}} \\ \midrule
\multicolumn{6}{l}{\textit{\textbf{N-Dimensional Efficiency ($D \in [3, 100]$)}}} \\ 
Feature Time (s) & 0.221 & 0.066 & 0.276 & \multicolumn{2}{c}{\textit{-- Failed --}} \\
Inference Time (s) & 0.109 & 0.066 & 0.276 & \multicolumn{2}{c}{\textit{-- Failed --}} \\
Total Time (s) & 0.330 & 0.132 & 0.552 & \multicolumn{2}{c}{\textit{-- Failed --}} \\
Speedup (\%) & -94.71 & 67.61 & 11.60 & \multicolumn{2}{c}{\textit{-- Failed --}} \\ \bottomrule
\end{tabular}%
}
\end{table}

\james{Lumping together dimensions 3-100 is a little much. Let's focus on smaller dimensions e.g., 3.  We can also acknowledge in higher dimensions, that feature generation dominates the time relative to, e.g., computing optimal solutions (since solution methods don't care about actual dimension). Also, explain speedup \% somewhere. Failed $\gets$ N/A or not predictable.  }

The breakdown of computational latency reveals that GART maintains a robust profile even in high dimensions. While the total prediction time increases to 0.330s in the N-Dimensional set, this remains strictly competitive compared to constructive heuristics like the Hilbert Curve (0.552s), which suffers from $O(D)$ overhead in the curve mapping process. Notably, the stability of GART ($\sigma=1.17$) drastically outperforms the benchmarks, ensuring reliable estimates regardless of dimensional scaling.




\section{Discussion} \label{sec:dsicussion}

The experimental results highlight a critical dichotomy in tour length estimation: the trade-off between geometric fidelity and computational overhead. Our analysis reveals that the utility of the Geometrically Assisted Regression Tree (GART) is intrinsically linked to the scale of the problem instance ($N$) and the dimensionality of the solution space ($D$).

\subsection{The Computational Bottleneck of Small-N Instances}
A fundamental limitation of graph-theoretic estimators, including GART, is their reliance on the distance matrix. Calculating the pairwise Euclidean distances for $N$ points requires $O(N^2 \cdot D)$ operations, followed by the $O(N^2)$ complexity of constructing the Minimum Spanning Tree (MST). For small instances ($N < 50$), this overhead is non-trivial relative to the time required to simply solve the TSP to optimality using efficient dynamic programming (e.g., Held-Karp) or branch-and-cut solvers like Concorde. In these "low-N" regimes, the computational cost of feature extraction for GART approaches the runtime of exact solvers, rendering the estimation step redundant. Therefore, GART is most effectively deployed as a filter for moderate-to-large scale problems ($N > 100$) where the exponential complexity of exact solvers makes real-time solution generation intractable, yet the quadratic cost of the MST remains negligible (milliseconds).

\subsection{Generalization in High-Dimensional Spaces}
The most significant contribution of GART is its robustness in high-dimensional hyperspaces ($D \gg 3$). Traditional asymptotic estimators, such as the BHH theorem \citep{bhh1959} and Chien's model \citep{chien1992operational}, rely on geometric priors—specifically the volume of the Convex Hull—to bound the feasible region. As demonstrated in our experiments, these methods suffer catastrophic failure modes as dimensions increase due to the "curse of dimensionality," where the Convex Hull calculation becomes computationally prohibitive and the concept of "volume" loses its descriptive power for point dispersion. 

GART circumvents this failure by discarding coordinate-dependent volume metrics in favor of topological invariants derived from the MST (e.g., dominance ratios, spectral diameter). This allows the model to generalize effectively to $D=100$ and beyond without retraining, providing high-fidelity estimates in scenarios where geometric intuition breaks down. This capability is particularly relevant for abstract logistics problems, such as DNA sequencing or vector quantization, which are often modeled as TSPs in high-dimensional feature spaces.

\subsection{Comparison with Deep Learning Approaches}
It is instructive to contrast GART with recent advancements in Deep Learning for combinatorial optimization. For instance, \citet{canturk2024scalable} recently proposed a Graph Neural Network (GNN) framework for scalable primal heuristics. While their approach demonstrates the power of deep representations in learning complex branching policies, it entails significant architectural complexity, requiring extensive training on GPU infrastructure and opaque latent feature extractions. In contrast, GART relies on a "shallow" Gradient Boosted Decision Tree architecture that utilizes a compact, interpretable set of topological features. While GNN-based models like those of \citet{canturk2024scalable} may offer marginal gains in solution construction, GART provides a lightweight, CPU-efficient alternative for pure \textit{estimation} tasks, effectively bridging the gap between naive theoretical bounds and heavy deep-learning models.

\section{Conclusion}\label{sec:conclusion}

This work presented the Geometrically Assisted Regression Tree (GART), a dimension-agnostic framework for estimating TSP tour lengths. By decoupling the estimation problem into a deterministic topological component (the MST) and a statistical correction factor ($\alpha$), GART overcomes the severe limitations of classical geometric estimators in high-dimensional and anisotropic spaces.

Our extensive evaluation demonstrates that while GART incurs an efficiency penalty for trivial instance sizes ($N < 50$) due to distance matrix calculations, it achieves optimal performance for large-scale operations ($N \ge 100$), delivering estimates with $<1\%$ error in millisecond-scale runtimes. Crucially, GART proves immune to the "curse of dimensionality" that renders Convex Hull-based models intractable in hyperspaces ($D=100$), validating its utility for complex, non-Euclidean routing problems.

Ultimately, GART occupies a unique niche in the operations research toolkit. It offers a computationally inexpensive yet statistically robust alternative to both asymptotic bounds and complex Graph Neural Networks, enabling real-time decision support for dynamic logistics systems where speed and accuracy are paramount.

%Bibliography
\bibliographystyle{plainnat}  
\bibliography{references}  


\end{document}
